\scanpage{96}
\diagram{\begin{tikzcd}[row sep=small,column sep=small]
0&0&0&\\
0&0&A^{2\,2}&\text{---}\\
0&A^{1\,1}&A^{2\,1}\arrow[u]&\text{---}\\
k&A^{1\,0}\arrow[u]&A^{2\,0}\arrow[u]&\cdots
\end{tikzcd}}

Soit $A^{**}\in\operatorname{Ob}\mathrm{DR}_{\unsure{\mathrm{sp}}}(X,k)$

$T\in\Gamma_X(A^{1,0})$ définit
$\begin{array}{rcl}k&\longrightarrow&A^{1,0}=\Phi,\\\lambda&\longmapsto&\lambda T\end{array}$,
\unsure{on a suite} exacte
\diagram{\begin{tikzcd}
0\arrow[r,dash]&k\arrow[r]&\Phi\arrow[r]\arrow[d,equal]&\Psi\arrow[r]\arrow[d,dash]&0\\
&&A^{1\,0}&A^{\unsure{2\,0}}
\end{tikzcd}}
Comme $\Psi$ est $k$-plat, $\Phi$ \unsure{aussi} et on sait que
\[
\mathrm{DR}(\Phi,T)=\underbrace{\Gamma^*(\Phi)}_{\text{deg}\atop\text{compl.}}\otimes_k\underbrace{\bigwedge^*\Psi}_{\text{deg ext.}}
\]
\unsure{muni} des \deleted{\unsure{str.}} \inserted{\unsure{structures}} habituelles est une $k$-alg. de De Rham \unsure{spéciale} et on a un $k$-hom. canonique de telles
\[
\mathrm{DR}(\Phi,T)\longrightarrow A^{**}
\]
\deleted{donc la cat. de}

Un hom $A^{**}\unsure{-|\longrightarrow}A'{}^{**}$ définit diagramme commutatif
\diagram{\begin{tikzcd}
\mathrm{DR}(\Phi,T)\arrow[r]\arrow[d]&A^{**}\arrow[d]\\
\mathrm{DR}(\Phi',T')\arrow[r]&A'{}^{**}
\end{tikzcd}}
Cela \unsure{implique} aisément que la cat. dérivée $\mathrm{DR}_{\unsure{\mathrm{sp}}}(X,k)_0$ est \unsure{équivalente} canoniquement à la \deleted{\unsure{sous-}}catégorie \inserted{dérivée} \unsure{déduite} des algèbres de la forme $\mathrm{DR}(\Phi,T)$, qui est aussi celle des suites exactes
\diagram{\begin{tikzcd}
0\arrow[r]&k\arrow[r]&\Phi\arrow[r]&\Psi\arrow[r]&0
\end{tikzcd}}
avec $\Psi$ plat sur $k$.

\margin{95}
